\section{Methodology}

The implemented pipeline consists of five stages: data preparation, model training, Rashomon-set selection, prediction-disagreement analysis, and explanation-variability analysis. 

\subsection{Data Preparation}

The current implementation uses \gls{NSLKDD} as the default intrusion-detection dataset. The dataset is downloaded automatically on the first run and stored locally under \texttt{data/raw/nsl-kdd/}. The task is configured as binary classification by default, mapping all attack labels to one attack class and retaining normal traffic as the second class. A multiclass configuration is supported in the data loader but is not the focus of the current draft.

Categorical features, namely protocol type, service, and flag, are transformed with one-hot encoding. Numeric features are standardized with a standard scaler. The preprocessing pipeline is fitted on the training split and then applied to validation and test data. The fitted preprocessor and feature names are stored in the run directory to keep later explanation steps aligned with the trained models.

\subsection{Model Family}

Each classifier is a feed-forward multilayer perceptron implemented in PyTorch. The architecture is configurable through the experiment configuration file. Hidden layers use linear transformations, ReLU activations, batch normalization, and dropout, followed by a final linear classification layer. Models are trained with AdamW and cross-entropy loss. Early stopping is based on validation loss.

\subsection{Generating a Candidate Model Set}

To induce model multiplicity, the pipeline trains multiple models on the same prediction task while varying random seeds and, optionally, the fraction of the training set. For subset experiments, examples are sampled with a stratified strategy so that class balance is approximately preserved. Each trained model is evaluated on the same test split, and the pipeline stores per-model metrics as well as per-sample predictions and predicted class probabilities.

\subsection{Rashomon-Set Selection}

After training, models are filtered into a Rashomon set. Given a metric such as accuracy and a tolerance value, the best observed score is identified and all models with score at least the best score minus the tolerance are retained. The tolerance therefore controls how broadly similarly good models are included in the analysis.

\subsection{Predictive Multiplicity}

Predictive multiplicity is measured at two levels. At the model-pair level, the pipeline computes pairwise disagreement rates between all models in the Rashomon set. At the sample level, predictions are pivoted into a model-by-sample vote matrix. For each sample, the implementation computes the majority prediction, vote entropy, the number of unique predictions, the majority fraction, and the conflict ratio. The conflict ratio is defined as one minus the majority fraction and is positive whenever at least one Rashomon model disagrees with the majority vote.

\subsection{Explanation Variability}

The pipeline computes \gls{SHAP} explanations for each model in the Rashomon set. A random subset of training samples is used as background data, and a configurable subset of test samples is explained. When enabled, the \texttt{--only-conflicts} option restricts the explanation set to samples with positive conflict ratio, matching the conflict-focused analysis used in the prior BA-style experiments.

For each model, class, sample, and feature, the \gls{SHAP} value is stored in long format. The implementation also computes mean absolute \gls{SHAP} values per feature and model, which are used to compare feature rankings and explanation distances. Pairwise explanation disagreement is summarized using cosine distance between mean-absolute-\gls{SHAP} vectors and Jaccard overlap between top-\(k\) feature sets.

\subsection{Variability Metrics}

Our analysis aggregates \gls{SHAP} values across models for each sample, class, and feature. It computes the minimum, maximum, mean, standard deviation, variance, and range of \gls{SHAP} values. In addition, sign instability is measured as the smaller of the positive and negative sign fractions across models. This captures whether similarly accurate models agree on the direction of a feature's contribution.

Finally, the implementation relates explanation variability to predictive multiplicity by computing Spearman correlations between conflict ratio and \gls{SHAP} value range, \gls{SHAP} value variance, and sign instability for each feature and class. The plotting module produces publication-oriented figures for metric variability, predictive disagreement, \gls{SHAP} feature importance, explanation disagreement, sign-instability heatmaps, correlation matrices, and per-feature scatter plots.
